\documentclass{article}
\usepackage[T1]{fontenc}
\usepackage{lmodern}
\usepackage{graphicx}
\usepackage{amsmath,amssymb}
\usepackage[a4paper,margin=2cm]{geometry}
\usepackage[absolute,overlay]{textpos}
\setlength{\TPHorizModule}{1mm}
\setlength{\TPVertModule}{1mm}
\usepackage[indonesian]{babel}
\usepackage{hyperref}
\setlength{\emergencystretch}{2em}
% unit: Halaman 14

\begin{document}

% === Halaman 14 (python/native) ===

14

Contoh pengiriman pesan 2:

• Misalkan Bob akan mengirim plainteks  M = ‘HELLO ALICE’ kepada Alice

• Dengan memisalkan A = 00, B = 01, …, Z = 25, maka plainteks dikodekan ke dalam 

integer  (spasi diabaikan) menjadi 

m = 07041111140011080204

 Nilai m ini sangat besar untuk dipangkatkan dan tidak berada di dalam selang         

[0, 3337-1], oleh karena itu pecah m menjadi blok-blok pesan yang lebih kecil, 

misalnya blok 4 digit:


m1 = 0704

m4 = 1108

m2 = 1111

m5 = 0204

m3 = 1400

   (Perhatikan, mi masih valid karena terletak di dalam selang [0, 3337 – 1] )

\end{document}
